\section{Exact-risk benchmark and centered comparator}
\label{app:known-risk}

The controlled experiment uses a Bernoulli variable $Z$ with
$\Pr(Z=1)=p$, target $Y=2Z-1$, and a two-coordinate simplex map
\[
 \Phi_a(0)=\left((1+a)/2,(1-a)/2\right),\qquad
 \Phi_a(1)=\left((1-a)/2,(1+a)/2\right),\qquad 0\leq a\leq1.
\]
The representation distance is $\sqrt2a$.  Write
$c=(2-\kappa\sqrt2a)_+$.  Proposition~\ref{prop:two-atom} gives the
exact minimum MSE $R^*=p(1-p)c^2$.  The target values, representation
family, and support are specified by the benchmark design.  They are
not inferred using calibration outcomes.

An independent training sample yields an empirical frequency $\widehat p$.
The chosen constrained empirical minimizer predicts
\begin{equation}
 \widehat f(0)=-1+\widehat p c,
 \qquad
 \widehat f(1)=1-(1-\widehat p)c.
 \label{eq:two-atom-fit}
\end{equation}
These values are in $[-1,1]$ and separated by
$2-c=\min\{2,\kappa\sqrt2a\}$, so an affine decoder realizes them
within the budget.  If one atom is absent from training, the formula
fixes one feasible empirical-risk minimizer using the benchmark's
predeclared support; this is a tie convention.  The fitted population
risk can be computed without a test-set estimate:
\begin{equation}
 R(\widehat f)
 =c^2\{(1-p)\widehat p^2+p(1-\widehat p)^2\}
 =R^*+c^2(\widehat p-p)^2.
 \label{eq:two-atom-fitted-risk}
\end{equation}
The pair witness equals $c^2/2$ on unequal calibration atoms and zero
otherwise; its expectation is exactly $R^*$.  This verifies coverage
against known population risks, rather than treating a finite test
estimate as ground truth.

For comparison, use a single reference class and the raw-target
Euclidean version of the earlier centered certificate.  If $\bar z$
is the calibration frequency in $n$ independent observations, its
empirical centered output and input variances are
\[
 \widehat s=4\bar z(1-\bar z),\qquad
 \widehat t=2a^2\bar z(1-\bar z).
\]
For a failure probability $\eta$, define
\begin{align}
 r&=\log(2/\eta)/n,\nonumber\\
 V_\eta&=(\widehat s-4\sqrt{r/2})_+,\nonumber\\
 T_\eta&=\min\{1,\widehat t+\sqrt{2r}+1/n\},\nonumber\\
 L^{\rm ctr}_\eta&=(\sqrt{V_\eta}-\kappa\sqrt{T_\eta})_+^2.
 \label{eq:centered-comparator}
\end{align}
Then $L^{\rm ctr}_\eta\leq R^*$ with probability at least $1-\eta$.
For completeness, the output sample variance has expectation at most
the population variance and changes by at most $4/n$ when one
observation changes.  Bounded differences therefore gives the lower
correction in $V_\eta$ at failure probability $\eta/2$.  The input
sample variance changes by at most $2/n$; its downward bias is
$\tau^2/n\leq1/n$, since a simplex-valued vector has total variance
at most one.  Bounded differences and the bias correction give
$T_\eta\geq\tau^2$ at failure probability $\eta/2$.  The union bound
and the centered population inequality prove the claim.  The constant
one here is the conservative general simplex-variance bound, not the
sharpest possible two-coordinate bound.

In each trial the fitted upper bound, the pair lower bound, and the
centered lower bound are each assigned failure probability $\delta/3$.
The two internal tails in \eqref{eq:centered-comparator} together use
one such allocation.  Thus the three reported directions have a joint
confidence guarantee within a trial.  There is no asserted global
$1-\delta$ event across all Monte Carlo trials, all grid conditions,
or all separately reported experiments.
